\documentclass[11pt]{article}
\usepackage{amsmath,amssymb,amsthm}
\usepackage[margin=1in]{geometry}
\usepackage{booktabs}

\newtheorem{theorem}{Theorem}
\newtheorem{lemma}[theorem]{Lemma}

\title{Problem 383: Divisibility Comparison between Factorials}
\author{Project Euler Solutions}
\date{}

\begin{document}
\maketitle

\section{Problem Statement}

Let $f_5(n)$ denote the largest integer $k$ such that $5^k \mid n$. Define $D(N) = \sum_{i=2}^{N} f_5\!\bigl(\binom{2i}{i}\bigr)$. Find $D(10^{12})$.

\section{Mathematical Foundation}

\begin{theorem}[Legendre's Formula]
For a prime $p$ and positive integer $n$:
\[
v_p(n!) = \sum_{j=1}^{\infty} \left\lfloor \frac{n}{p^j} \right\rfloor = \frac{n - s_p(n)}{p - 1},
\]
where $s_p(n)$ is the sum of the base-$p$ digits of $n$.
\end{theorem}

\begin{proof}
The number of multiples of $p^j$ in $\{1, \ldots, n\}$ is $\lfloor n/p^j \rfloor$. Summing over all $j$ counts each factor of $p$ in $n!$ exactly once. For the second equality, write $n = \sum d_j p^j$ and telescope.
\end{proof}

\begin{theorem}[Kummer's Theorem]
For a prime $p$ and $m \ge n \ge 0$:
\[
v_p\binom{m}{n} = c_p(n,\, m-n),
\]
where $c_p(n, m-n)$ is the number of carries when adding $n$ and $m - n$ in base $p$.
\end{theorem}

\begin{proof}
From Legendre's formula:
\[
v_p\binom{m}{n} = \frac{s_p(n) + s_p(m-n) - s_p(m)}{p-1}.
\]
The numerator equals $(p-1)$ times the number of carries in the base-$p$ addition $n + (m-n) = m$.
\end{proof}

\begin{lemma}[Central Binomial 5-adic Valuation]
For $p = 5$:
\[
v_5\binom{2i}{i} = \text{number of carries when adding } i + i \text{ in base } 5.
\]
\end{lemma}

\begin{proof}
Direct application of Kummer's theorem with $m = 2i$, $n = i$.
\end{proof}

\begin{lemma}[Carry Propagation]
When doubling $i$ in base $5$, a carry at position $j$ occurs iff $2d_j + c_j \ge 5$, where $d_j$ is the $j$-th base-5 digit and $c_j \in \{0,1\}$ is the incoming carry ($c_0 = 0$). The outgoing carry is $c_{j+1} = \lfloor(2d_j + c_j)/5\rfloor \in \{0,1\}$.
\end{lemma}

\begin{proof}
Since $d_j \in \{0,\ldots,4\}$ and $c_j \in \{0,1\}$, we have $2d_j + c_j \le 9 < 2 \cdot 5$, so the carry is at most~1.
\end{proof}

\section{Algorithm}

The sum $D(N)$ counts total carries when doubling each $i \in [2, N]$ in base~5. This admits a digit DP over the base-5 representation of $N$.

\begin{verbatim}
function D(N):
    digits = base5_digits(N)
    DP state: (position, carry_in, tight)
    DP value: (count, total_carries)
    For each digit position, expand all (digit, carry) transitions
    Return total carries summed over terminal states
\end{verbatim}

\section{Complexity Analysis}

\begin{itemize}
\item \textbf{Time:} $O(\log_5 N)$. The DP has $O(\log_5 N)$ positions with $O(1)$ states each (at most $5 \times 2 \times 2 = 20$ transitions per position).
\item \textbf{Space:} $O(\log N)$.
\end{itemize}

\section{Answer}

\[
\boxed{22173624649806}
\]

\end{document}
